\documentclass[11pt,a4paper]{article}

\usepackage[T1]{fontenc}
\usepackage[utf8]{inputenc}
\usepackage[a4paper,margin=27mm]{geometry}
\usepackage{amsmath,amssymb,amsthm,mathtools}
\usepackage{newpxtext,newpxmath}
\usepackage{microtype}
\usepackage{booktabs,tabularx,array}
\usepackage{xcolor}
\usepackage{titlesec}
\usepackage[authoryear,round]{natbib}
\usepackage[colorlinks=true,linkcolor=blue!45!black,
  citecolor=red!45!black,urlcolor=blue!45!black]{hyperref}
\usepackage[nameinlink,noabbrev]{cleveref}

\definecolor{otblue}{HTML}{1F4E79}
\titleformat{\section}{\normalfont\Large\bfseries\color{otblue}}
  {\thesection}{.7em}{}
\titleformat{\subsection}{\normalfont\large\bfseries\color{otblue!80!black}}
  {\thesubsection}{.65em}{}
\titleformat{\paragraph}[runin]
  {\normalfont\bfseries\color{otblue!75!black}}{}{0pt}{}
\titlespacing*{\section}{0pt}{2.7ex plus .7ex minus .2ex}{1.1ex}
\titlespacing*{\subsection}{0pt}{2ex plus .5ex minus .2ex}{.75ex}
\titlespacing*{\paragraph}{0pt}{1.4ex plus .3ex minus .1ex}{.55em}

\newtheorem{theorem}{Theorem}[section]
\newtheorem{proposition}[theorem]{Proposition}
\newtheorem{lemma}[theorem]{Lemma}
\newtheorem{corollary}[theorem]{Corollary}
\theoremstyle{definition}
\newtheorem{definition}[theorem]{Definition}
\newtheorem{example}[theorem]{Example}
\theoremstyle{remark}
\newtheorem{remark}[theorem]{Remark}

\newcommand{\R}{\mathbb R}
\newcommand{\Pcal}{\mathcal P}
\newcommand{\Mcal}{\mathcal M}
\newcommand{\Tcal}{\mathcal T}
\newcommand{\Qcal}{\mathcal Q}
\newcommand{\Bcal}{\mathcal B}
\newcommand{\Piopt}{\Pi}
\newcommand{\dd}{\,\mathrm d}
\newcommand{\ones}{\mathbf 1}
\newcommand{\abs}[1]{\left\lvert #1\right\rvert}

\hypersetup{
  pdftitle={Conditional Concavity of p-Gromov--Wasserstein Energies},
  pdfauthor={Gabriel Peyre},
  pdfsubject={The exceptional quadratic exponent, counterexamples, and submodular ultrametric extensions},
  pdfkeywords={Gromov--Wasserstein, conditional concavity, negative type, ultrametric, submodularity, layer-cake formula}
}

\title{\textbf{Conditional Concavity of\\
\(p\)-Gromov--Wasserstein Energies}\\
\large Why the quadratic exponent is exceptional, and a valid all-exponent extension}
\author{Gabriel Peyr\'e}
\date{July 2026}

\begin{document}
\maketitle

\begin{abstract}
The usual conditional-concavity theorem for Gromov--Wasserstein transport
uses the squared distortion
\(\lvert d_{\mathcal X}-d_{\mathcal Y}\rvert^2\) and same-sign
conditionally definite base dissimilarities.  It is tempting to replace the
outer exponent \(2\) by an arbitrary \(p\geq1\).  This note shows that such a
replacement is false under the original assumptions: for every
\(p\neq2\), explicit four-point CND metric spaces have a positive feasible
curvature direction.  Thus \(p=2\) is genuinely exceptional, not a
technical limitation of the standard proof.

An exponent-uniform theorem nevertheless holds under a stronger geometric
hypothesis.  If all threshold kernels
\(\mathbf 1_{\{d\geq r\}}\) of the two dissimilarities are conditionally
negative definite, then every submodular scalar distortion with a nonnegative
mixed measure---in particular \(\lvert d_{\mathcal X}-d_{\mathcal Y}\rvert^p\)
for \(p\geq1\)---has nonpositive curvature on zero-marginal directions.  For
genuine metrics, this threshold condition is equivalent to ultrametricity.
A single layer-cake identity unifies this result with the classical quadratic
theory.  At \(p=2\), its scale weight is constant and integration recovers the
usual CND tensor factorization; on ultrametric spaces, the same integrand is
nonnegative at every pair of scales, which covers all \(p\geq1\).  The finite
result extends to arbitrary probability measures on compact ultrametric
spaces.  An extended catalogue connects this hypothesis to finite
hierarchies, symbolic sequences, valued fields, profinite groups, rooted
graphs and random tree boundaries.  In the finite uniform case, concavity
also guarantees a permutation optimizer.
\end{abstract}

\section{Curvature along the coupling constraints}
\label{sec:setup}

The GW objective is quadratic in the coupling for every fixed distortion
kernel.  Its convexity or concavity is therefore decided exactly by the sign
of one quadratic form on directions that preserve both marginals.

Let \(\mathcal X\) and \(\mathcal Y\) be compact spaces, equipped with bounded
symmetric dissimilarities
\(d_{\mathcal X}:\mathcal X^2\to\R_+\) and
\(d_{\mathcal Y}:\mathcal Y^2\to\R_+\).  For \(p\geq1\), define
\begin{equation}
  L_p\bigl((x,y),(x',y')\bigr)
  :=
  \abs{d_{\mathcal X}(x,x')-d_{\mathcal Y}(y,y')}^p.
  \label{eq:lp-kernel}
\end{equation}
For finite signed measures \(\xi,\eta\) on
\(\mathcal Z:=\mathcal X\times\mathcal Y\), set
\begin{equation}
  \Bcal_p(\xi,\eta)
  :=\iint_{\mathcal Z^2}L_p(z,z')\dd\xi(z)\dd\eta(z'),
  \qquad
  \Qcal_p(\xi):=\Bcal_p(\xi,\xi).
  \label{eq:quadratic-form}
\end{equation}
The \(p\)-GW problem minimizes \(\Qcal_p(\pi)\) over
\(\pi\in\Piopt(\alpha,\beta)\).  Its feasible direction space is
\begin{equation}
  \Tcal_{\mathcal X,\mathcal Y}
  :=
  \left\{\xi\in\Mcal(\mathcal X\times\mathcal Y):
  (P_{\mathcal X})_\#\xi=0,\ (P_{\mathcal Y})_\#\xi=0\right\}.
  \label{eq:tangent-space}
\end{equation}

\begin{proposition}[Exact feasible-curvature criterion]
\label{prop:curvature-criterion}
For \(\pi_0,\pi_1\in\Piopt(\alpha,\beta)\), write
\(\xi=\pi_1-\pi_0\).  Then
\begin{equation}
  \frac{\mathrm d^2}{\mathrm dt^2}
  \Qcal_p\bigl((1-t)\pi_0+t\pi_1\bigr)
  =2\Qcal_p(\xi).
  \label{eq:segment-curvature}
\end{equation}
Consequently, \(\Qcal_p\) is concave on
\(\Piopt(\alpha,\beta)\) whenever
\(\Qcal_p(\xi)\leq0\) for every
\(\xi\in\Tcal_{\mathcal X,\mathcal Y}\).  For finite spaces and strictly
positive marginals, this condition is also necessary.
\end{proposition}

\begin{proof}
Expanding the quadratic form gives
\[
  \Qcal_p(\pi_0+t\xi)
  =\Qcal_p(\pi_0)+2t\Bcal_p(\pi_0,\xi)+t^2\Qcal_p(\xi),
\]
which proves \eqref{eq:segment-curvature}.  Every difference of feasible
couplings has zero marginals, so the sign condition implies concavity on each
feasible segment.  In the finite positive-marginal case, the product coupling
lies in the relative interior of the transport polytope.  Hence every
zero-row- and zero-column-sum matrix is, after multiplication by a sufficiently
small scalar of either sign, a feasible direction through that coupling.
Concavity therefore forces its quadratic value to be nonpositive.
\end{proof}

\section{One multiscale identity, two concavity certificates}
\label{sec:master-decomposition}

The quadratic and ultrametric theories are two levels of the same multiscale
argument.  The common object records the curvature contributed by every pair
of distance thresholds.  The quadratic loss averages those contributions
with a constant weight, whereas a general power uses a non-flat scale measure
that is singular on the diagonal when \(p=1\), and therefore requires finer
control.

For a dissimilarity \(d\) and a threshold \(r>0\), define
\begin{equation}
  H_r^d(x,x'):=\mathbf 1_{\{d(x,x')\geq r\}}.
  \label{eq:threshold-kernel}
\end{equation}
For \(\xi\in\Tcal_{\mathcal X,\mathcal Y}\), its two-scale interaction is
\begin{equation}
  I_\xi(r,s)
  :=\iint
  H_r^{d_{\mathcal X}}(x,x')H_s^{d_{\mathcal Y}}(y,y')
  \dd\xi(x,y)\dd\xi(x',y').
  \label{eq:threshold-product}
\end{equation}

The same decomposition applies beyond power losses.  Let
\(D_{\mathcal X}\) and \(D_{\mathcal Y}\) bound the two dissimilarities.

\begin{definition}[Submodular mixed-measure distortion]
\label{def:mixed-measure-distortion}
A continuous function
\(\ell:[0,D_{\mathcal X}]\times[0,D_{\mathcal Y}]\to\R\) is a
\emph{submodular mixed-measure distortion} if there is a finite positive Borel
measure \(\omega_\ell\) on
\((0,D_{\mathcal X}]\times(0,D_{\mathcal Y}]\) such that
\begin{equation}
  \ell(a,b)
  =\ell(a,0)+\ell(0,b)-\ell(0,0)
  -\omega_\ell\bigl((0,a]\times(0,b]\bigr).
  \label{eq:mixed-measure-representation}
\end{equation}
\end{definition}

This is the measure-valued form of submodularity.  If \(\ell\) is twice
continuously differentiable, the condition
\(\partial_{ab}^2\ell\leq0\) gives
\(\dd\omega_\ell(r,s)=-\partial_{ab}^2\ell(r,s)\dd r\dd s\).
The measure formulation also accommodates nonsmooth distortions concentrated
on lower-dimensional sets.

\begin{proposition}[Master multiscale curvature identity]
\label{prop:master-multiscale}
For a finite signed measure
\(\xi\in\Tcal_{\mathcal X,\mathcal Y}\), define
\begin{equation}
  \Qcal_\ell(\xi)
  :=\iint
  \ell\bigl(d_{\mathcal X}(x,x'),d_{\mathcal Y}(y,y')\bigr)
  \dd\xi(x,y)\dd\xi(x',y').
  \label{eq:general-distortion-form}
\end{equation}
If \(\ell\) satisfies \cref{def:mixed-measure-distortion}, then
\begin{equation}
  \Qcal_\ell(\xi)
  =-\int I_\xi(r,s)\dd\omega_\ell(r,s).
  \label{eq:mixed-measure-curvature}
\end{equation}
In particular, for \(\ell_p(a,b)=\abs{a-b}^p\),
\begin{align}
  \Qcal_1(\xi)
  &=-2\int_0^\infty I_\xi(r,r)\dd r,
  \label{eq:p1-layer-cake}\\
  \Qcal_p(\xi)
  &=-p(p-1)\int_0^\infty\int_0^\infty
  \abs{r-s}^{p-2}I_\xi(r,s)\dd r\dd s,
  \qquad p>1.
  \label{eq:p-layer-cake-curvature}
\end{align}
\end{proposition}

\begin{proof}
Insert \eqref{eq:mixed-measure-representation} into
\eqref{eq:general-distortion-form}.  Terms depending only on
\(d_{\mathcal X}\), only on \(d_{\mathcal Y}\), or on neither vanish because
\(\xi\) has zero marginals.  Moreover,
\[
  \omega_\ell((0,a]\times(0,b])
  =\int\mathbf 1_{\{a\geq r\}}\mathbf 1_{\{b\geq s\}}
  \dd\omega_\ell(r,s).
\]
Fubini's theorem therefore gives \eqref{eq:mixed-measure-curvature}.

For \(p=1\), use
\(a+b-\abs{a-b}=2\int_0^\infty
\mathbf 1_{\{a\geq r\}}\mathbf 1_{\{b\geq r\}}\dd r\).
For \(p>1\), two integrations of the mixed derivative give
\[
  a^p+b^p-\abs{a-b}^p
  =p(p-1)\int_0^a\int_0^b\abs{r-s}^{p-2}\dd r\dd s.
\]
The kernel \(\abs{r-s}^{p-2}\) is locally integrable exactly for \(p>1\),
which proves both power formulas.
\end{proof}

\paragraph{The quadratic exponent integrates the scales.}
Setting \(p=2\) in \eqref{eq:p-layer-cake-curvature} makes the scale weight
equal to one.  Since
\(\int_0^\infty H_r^d(x,x')\dd r=d(x,x')\), Fubini gives
\begin{equation}
  \Qcal_2(\xi)
  =-2\int_0^\infty\int_0^\infty I_\xi(r,s)\dd r\dd s
  =-2\iint
  d_{\mathcal X}(x,x')d_{\mathcal Y}(y,y')
  \dd\xi(x,y)\dd\xi(x',y').
  \label{eq:p2-factorization}
\end{equation}
Thus the familiar quadratic factorization is not a separate theory: it is the
flat-weight member of the multiscale identity.

\begin{corollary}[Integrated quadratic certificate]
\label{cor:quadratic-integrated-certificate}
If \(d_{\mathcal X}\) and \(d_{\mathcal Y}\) are both conditionally negative
definite (CND), or both conditionally positive definite (CPD), then
\(\Qcal_2(\xi)\leq0\) for every
\(\xi\in\Tcal_{\mathcal X,\mathcal Y}\).  Hence squared GW is concave on each
transport polytope.
\end{corollary}

\begin{proof}
Center each dissimilarity at an arbitrary base point.  Additive one-variable
terms disappear against the zero marginals of \(\xi\).  In the CND case, the
negatives of the two centered kernels are positive semidefinite; in the CPD
case, the centered kernels themselves are positive semidefinite.  Their
Kronecker-product quadratic form is therefore nonnegative in both cases.
Equation \eqref{eq:p2-factorization} gives the result.
\end{proof}

This is the mechanism behind conditional concavity in graph matching and
squared GW \citep{MaronLipman2018,SejourneVialardPeyre2021}.  Crucially, only
the \emph{integrated} kernels \(d_{\mathcal X}\) and
\(d_{\mathcal Y}\) need to have the appropriate sign; the individual
threshold interactions \(I_\xi(r,s)\) may change sign.

The word \emph{squared} here refers to the outer exponent in
\eqref{eq:lp-kernel}.  It should not be confused with an inner power in the
base dissimilarities.  For example, Schoenberg's theorem implies that
\(\lVert x-x'\rVert^q\) is CND for \(0<q\leq2\)
\citep{Schoenberg1938}; each such inner exponent is covered by
\eqref{eq:p2-factorization}, but the outer exponent remains exactly \(2\).
Recent Hilbert-embedding formulations likewise broaden the base costs while
retaining a squared outer discrepancy \citep{HouryFeydyVialard2026}.

\paragraph{Ultrametric geometry controls every scale separately.}
Formula \eqref{eq:mixed-measure-curvature} reveals a stronger sufficient
condition: if \(I_\xi(r,s)\geq0\) for every feasible direction \(\xi\) and
every pair of thresholds, then every submodular mixed-measure distortion is
conditionally concave.  Threshold-CND dissimilarities provide exactly this
pointwise certificate.  For genuine metrics, the next sections show that this
condition is equivalent to ultrametricity.  Consequently, the ultrametric
theory contains the quadratic result as its \(p=2\) member, but strengthens
its hypothesis in order to cover all the other scale measures as well.

\section{The CND hypothesis does not extend beyond \texorpdfstring{\(p=2\)}{p=2}}
\label{sec:counterexamples}

The collapse of the multiscale formula to the product factorization
\eqref{eq:p2-factorization} is specific to its constant \(p=2\) weight.  For
other exponents, \cref{prop:master-multiscale} retains scale-dependent
information whose sign is not controlled by ordinary CND.  The failure
already occurs for four-point path, tree and Hamming metrics, so it persists
even when the base kernels are genuine CND distances rather than arbitrary
CND dissimilarities.

\begin{theorem}[The original hypothesis is specific to \(p=2\)]
\label{thm:no-cnd-extension}
For every \(p\in[1,\infty)\setminus\{2\}\), there are four-point metric
spaces \((\mathcal X,d_{\mathcal X})\) and
\((\mathcal Y,d_{\mathcal Y})\) whose distance kernels are CND, and a
nonzero matrix \(\Xi\in\Tcal_{\mathcal X,\mathcal Y}\), such that
\(\Qcal_p(\Xi)>0\).  Hence same-sign conditional definiteness of genuine
metric distances does not imply conditional concavity for any outer exponent
other than \(2\).
\end{theorem}

\begin{proof}
We use only finite \(\ell^1\) and tree metrics.  They are CND: in one
dimension,
\[
  \abs{s-t}
  =\int_{\R}
  \abs{\mathbf 1_{\{s>r\}}-\mathbf 1_{\{t>r\}}}^2\dd r,
\]
so, for \(\sum_i a_i=0\), its quadratic form equals
\(-2\int_{\R}\abs{\sum_i a_i\mathbf 1_{\{s_i>r\}}}^2\dd r\leq0\).
Finite sums preserve conditional negative definiteness, proving the claim for
\(\ell^1\); tree metrics are sums of the same cut metrics.

First let \(1\leq p<2\).  Take the line metric on the labeled points
\((0,3,2,1)\) and the Hamming metric on the labeled square
\(((0,0),(1,1),(1,0),(0,1))\).  Their distance matrices are
\begin{equation}
  \mathrm D_{\mathcal X}
  =\begin{pmatrix}
    0&3&2&1\\3&0&1&2\\2&1&0&1\\1&2&1&0
  \end{pmatrix},
  \qquad
  \mathrm D_{\mathcal Y}
  =\begin{pmatrix}
    0&2&1&1\\2&0&1&1\\1&1&0&2\\1&1&2&0
  \end{pmatrix}.
  \label{eq:small-p-metrics}
\end{equation}
Set \(u=(1,-1,-1,1)^\top\), \(v=(1,1,-1,-1)^\top\) and
\(\Xi_{<}=uv^\top\).  Since both vectors sum to zero, \(\Xi_{<}\) has zero
row and column sums.  A direct collection of the distance differences in
\eqref{eq:small-p-metrics} gives
\begin{equation}
  \Qcal_p(\Xi_{<})
  =8\bigl(2^{p+1}-3^p+1\bigr)>0.
  \label{eq:small-p-counterexample}
\end{equation}
Indeed, \(\varphi(p)=(3/2)^p-2^{-p}\) is strictly increasing and
\(\varphi(2)=2\).  Hence \(p<2\) implies
\(3^p-1<2^{p+1}\), proving the strict inequality in
\eqref{eq:small-p-counterexample}.  This includes \(p=1\).

Now let \(p>2\).  Keep \(\mathrm D_{\mathcal Y}\) from
\eqref{eq:small-p-metrics}, but let \(\mathcal X\) be the unit-edge star with
three leaves and its center as the fourth point:
\begin{equation}
  \mathrm D_{\mathcal X}
  =\begin{pmatrix}
    0&2&2&1\\2&0&2&1\\2&2&0&1\\1&1&1&0
  \end{pmatrix}.
  \label{eq:large-p-metric}
\end{equation}
With \(\widetilde u=(1,0,0,-1)^\top\), use the feasible direction
\(\Xi_{>}=\widetilde u v^\top\).  Direct summation now yields
\begin{equation}
  \Qcal_p(\Xi_{>})=8\bigl(2^p-4\bigr)>0,
  \label{eq:large-p-counterexample}
\end{equation}
which is strictly positive exactly when \(p>2\).
\end{proof}

\begin{remark}[What the counterexamples rule out]
The theorem does not say that every \(p\neq2\) GW objective is nonconcave.  It
says that CND or CPD base dissimilarities alone no longer certify concavity.
Any valid extension must impose additional structure on the base geometry or
directly on the distortion kernel.  It also shows why replacing the tensor
factorization in the \(p=2\) proof by an appeal to Schoenberg's theorem would
be invalid: Schoenberg controls powers of a single distance kernel, whereas
\(\abs{d_{\mathcal X}-d_{\mathcal Y}}^p\) couples two distance values through
a nonlinear difference.
\end{remark}

\section{Pointwise threshold positivity and ultrametrics}
\label{sec:thresholds}

The master identity isolates the extra structure needed away from the
quadratic exponent: positivity must hold before the distance scales are
averaged.  Requiring conditional negative type at every threshold supplies a
clean geometric certificate.

\begin{definition}[Threshold negative type]
\label{def:threshold-cnd}
A dissimilarity is \emph{threshold-CND} if every kernel \(H_r^d\) defined in
\eqref{eq:threshold-kernel} is CND.  On an infinite space, this means that the
CND inequality holds on every finite subset.
\end{definition}

Threshold-CND implies CND because
\(d(x,x')=\int_0^\infty H_r^d(x,x')\dd r\).  For genuine metrics, the stronger
condition has an exact geometric meaning.

\begin{proposition}[Threshold negative type characterizes ultrametrics]
\label{prop:threshold-ultrametric}
A metric \(d\) is threshold-CND if and only if it is an ultrametric, namely
\begin{equation}
  d(x,z)\leq\max\{d(x,y),d(y,z)\}
  \qquad\text{for all }x,y,z.
  \label{eq:ultrametric-inequality}
\end{equation}
\end{proposition}

\begin{proof}
Suppose first that \(d\) is ultrametric.  For each \(r>0\), the relation
\(x\sim_r x'\) if \(d(x,x')<r\) is an equivalence relation.  On any finite
collection \((x_i)_i\), let \(\mathscr P_r\) denote the resulting partition.
For coefficients satisfying \(\sum_i a_i=0\),
\[
  \sum_{i,j}a_i a_jH_r^d(x_i,x_j)
  =-\sum_{C\in\mathscr P_r}\left(\sum_{i:x_i\in C}a_i\right)^2\leq0.
\]
Thus every \(H_r^d\) is CND.

Conversely, if \eqref{eq:ultrametric-inequality} fails, there are
\(x,y,z\) such that
\(d(x,z)>\max\{d(x,y),d(y,z)\}\).  Choose \(r\) strictly above the latter
maximum and no larger than \(d(x,z)\).  The threshold kernel vanishes on
\((x,y)\) and \((y,z)\), but equals one on \((x,z)\).  Its quadratic form at
the zero-sum coefficients \((1,-2,1)\) is therefore \(2>0\), contradicting
conditional negative definiteness.
\end{proof}

Thus threshold-CND is much stronger than ordinary negative type among metrics:
it selects precisely ultrametric geometry.  The line and tree metrics used in
\cref{thm:no-cnd-extension} are CND but not ultrametric.

\begin{theorem}[Submodular distortions under threshold negative type]
\label{thm:threshold-concavity}
Let \(\mathcal X\) and \(\mathcal Y\) be finite spaces whose dissimilarities are
threshold-CND, and let \(\ell\) satisfy
\cref{def:mixed-measure-distortion}.  Then
\(\Qcal_\ell(\Xi)\leq0\) for every
\(\Xi\in\Tcal_{\mathcal X,\mathcal Y}\).  Consequently, the associated
quadratic transport objective is concave on every transport polytope over the
two spaces.
\end{theorem}

\begin{proof}
Since both threshold kernels are CND, centering their negatives at arbitrary
base points produces positive semidefinite matrices.  All centering terms
disappear because \(\Xi\) has zero row and column sums.  Hence
\(I_\Xi(r,s)\) is the quadratic form of a Kronecker product of two positive
semidefinite matrices, and is nonnegative; the two minus signs cancel.
The master identity \eqref{eq:mixed-measure-curvature} now gives
\(\Qcal_\ell(\Xi)\leq0\).
\end{proof}

\begin{corollary}[All power distortions]
\label{cor:power-distortions}
Under the assumptions of \cref{thm:threshold-concavity}, the choice
\(\ell_p(a,b)=\abs{a-b}^p\) gives
\(\Qcal_p(\Xi)\leq0\) for every \(p\geq1\).  In particular, \(p=2\) is
included as the constant-weight specialization of
\eqref{eq:p-layer-cake-curvature}.
\end{corollary}

\begin{proof}
Apply \cref{thm:threshold-concavity} to the power-loss representations in
\cref{prop:master-multiscale}.
\end{proof}

\begin{corollary}[Sparse and permutation minimizers]
\label{cor:ultrametric-extreme-minimizer}
Let \(\mathcal X\) and \(\mathcal Y\) be finite ultrametric spaces with
strictly positive marginals supported on \(n\) and \(m\) points.  For every
\(p\geq1\), the \(p\)-GW problem admits an optimal coupling supported on at
most \(n+m-1\) pairs.  If \(n=m\) and both marginals are uniform, it admits a
permutation coupling as an optimizer.
\end{corollary}

\begin{proof}
By \cref{prop:threshold-ultrametric,cor:power-distortions}, the objective is
concave on the transport polytope.  Write any minimizer as a convex combination
of extreme points.  Concavity shows that its value is no smaller than the
corresponding convex combination of their values; optimality forces every
extreme point used with positive weight to be optimal.  An extreme transport
plan has at most \(n+m-1\) positive entries.  For equal uniform marginals, the
Birkhoff--von Neumann theorem identifies the extreme points with scaled
permutation matrices.
\end{proof}

\paragraph{Why this is the right strengthening.}
The two theories are now nested rather than disconnected.  Ordinary CND
controls the scale average
\(d=\int H_r^d\dd r\), which suffices when \(p=2\) because the weight in
\eqref{eq:p-layer-cake-curvature} is constant.  Threshold-CND controls
\(I_\xi(r,s)\) pointwise, so every nonnegative mixed measure may be used; for
metrics, this is precisely the ultrametric hypothesis.  The latter therefore
recovers the quadratic theorem and extends it to all \(p\geq1\).  It is
sufficient, but not claimed to be necessary for concavity at one fixed
exponent and one fixed pair of spaces: the exact finite criterion remains
\cref{prop:curvature-criterion}.

\section{A hierarchy of ultrametric examples}
\label{sec:ultrametric-examples}

Ultrametrics range from the constant distance on a finite set to spaces of
infinite trees, valued fields and locally convergent graphs.  Their common
mechanism is not an exotic triangle inequality but a nested family of
partitions: two points are close when they remain in the same cluster down to
a fine scale.  This section builds a catalogue from elementary examples to
more structural ones and records compactness, because compact ultrametric
spaces enter directly into \cref{thm:ultrametric-concavity}.

\paragraph{Nested partitions are the basic construction.}
The finite case already contains the complete geometry.  It also explains
why dendrograms and ultrametrics are equivalent representations of the same
hierarchy \citep{CarlssonMemoli2010}.

\begin{proposition}[Finite hierarchies and ultrametrics]
\label{prop:nested-partition-ultrametric}
Let \(\mathscr P_0,\ldots,\mathscr P_L\) be partitions of a finite set
\(\mathcal X\), where \(\mathscr P_0\) is the partition into singletons,
\(\mathscr P_L=\{\mathcal X\}\), and each \(\mathscr P_\ell\) refines
\(\mathscr P_{\ell+1}\).  Fix heights
\(0=h_0<h_1<\cdots<h_L\).  For \(x\neq x'\), set
\begin{equation}
  m(x,x'):=\min\{\ell:x,x'\text{ lie in the same block of }
  \mathscr P_\ell\},
  \qquad d(x,x'):=h_{m(x,x')},
  \label{eq:partition-ultrametric}
\end{equation}
and set \(d(x,x)=0\).  Then \(d\) is an ultrametric.  Conversely, every
finite ultrametric is obtained in this way.
\end{proposition}

\begin{proof}
If \(x,y\) and \(y,z\) belong to common blocks by level
\(\ell=\max\{m(x,y),m(y,z)\}\), those blocks intersect at \(y\) and hence are
the same block.  Thus \(m(x,z)\leq\ell\), which gives
\(d(x,z)\leq\max\{d(x,y),d(y,z)\}\).

Conversely, let \(0=h_0<h_1<\cdots<h_L\) be the distinct values of a finite
ultrametric.  At level \(\ell\), declare \(x\sim_\ell x'\) when
\(d(x,x')\leq h_\ell\).  The strong triangle inequality makes this relation
transitive, so its classes form a partition \(\mathscr P_\ell\).  These
partitions are nested, and the first level at which \(x\) and \(x'\) merge
has height exactly \(d(x,x')\).
\end{proof}

\begin{example}[The constant discrete ultrametric]
\label{ex:constant-ultrametric}
On any set \(\mathcal X\), the metric
\[
  d(x,x')=R\,\mathbf 1_{\{x\neq x'\}},\qquad R>0,
\]
is ultrametric.  Every triangle is either degenerate or equilateral.  On a
finite set, this is \eqref{eq:partition-ultrametric} with only the singleton
partition and the one-block partition.  It is the simplest example, but also
the least informative one: there is only one nontrivial distance threshold.
\end{example}

\begin{example}[A two-scale block model]
\label{ex:two-scale-ultrametric}
Let \(0<r<R\), split \(\{1,2,3,4\}\) into the two blocks
\(\{1,2\}\) and \(\{3,4\}\), and use distance \(r\) inside a block and
\(R\) across blocks.  The distance matrix is
\[
  \begin{pmatrix}
    0&r&R&R\\
    r&0&R&R\\
    R&R&0&r\\
    R&R&r&0
  \end{pmatrix}.
\]
At thresholds up to \(r\), every point is isolated; between \(r\) and \(R\),
the two prescribed blocks are visible; above \(R\), everything has merged.
Replacing the two blocks by an arbitrary nested collection gives a finite
hierarchical block model.
\end{example}

\begin{example}[Dendrograms and weighted rooted trees]
\label{ex:dendrogram-ultrametric}
Let the points be the leaves of a rooted tree.  Give every vertex a height,
with leaves at height zero and heights strictly increasing toward the root.
If \(x\wedge x'\) denotes the lowest common ancestor of two leaves, then
\begin{equation}
  d(x,x')=h(x\wedge x')
  \label{eq:dendrogram-ultrametric}
\end{equation}
is ultrametric: among three leaves, the two largest merge heights coincide.
Conversely, \cref{prop:nested-partition-ultrametric} turns every finite
ultrametric into such a dendrogram.  This is the geometry produced by
hierarchical clustering and used for taxonomies and phylogenetic
representations.  It is important not to confuse \eqref{eq:dendrogram-ultrametric}
with the shortest-path metric on all vertices of a tree; the latter is
generally a tree metric, but not an ultrametric.
\end{example}

\paragraph{Prefix spaces turn hierarchy into coordinates.}
Words and sequences provide the canonical infinite-scale model: distance is
determined by the first coordinate at which two objects disagree.

\begin{example}[Finite words, Cantor space and Baire space]
\label{ex:prefix-ultrametric}
Let \(\mathcal A\) be an alphabet, let \(0<\vartheta<1\), and consider
sequences \(x=(x_k)_{k\geq1}\) in \(\mathcal A^{\mathbb N}\).  If
\(x\neq x'\), define the first disagreement
\(N(x,x')=\min\{k\geq1:x_k\neq x'_k\}\) and
\begin{equation}
  d_\vartheta(x,x')=\vartheta^{N(x,x')-1},
  \qquad d_\vartheta(x,x)=0.
  \label{eq:prefix-ultrametric}
\end{equation}
If two pairs agree through the first \(m\) symbols, then so does the remaining
pair; hence \eqref{eq:prefix-ultrametric} is ultrametric.  Its balls are
cylinder sets obtained by fixing an initial word.  Fixed-length words inherit
finite ultrametrics.  If \(\mathcal A\) is finite, the infinite sequence
space is compact and is usually called a Cantor space; if
\(\mathcal A=\mathbb N\), one obtains the complete but noncompact Baire
space.  Closed subshifts and spaces of admissible symbolic trajectories are
ultrametric subspaces of the same form.
\end{example}

\begin{example}[The ternary Cantor set]
\label{ex:cantor-set-ultrametric}
The coding map
\[
  \iota(\omega)=\sum_{k\geq1}2\omega_k3^{-k},
  \qquad \omega\in\{0,1\}^{\mathbb N},
\]
identifies binary sequences with the middle-third Cantor set.  Transporting
\eqref{eq:prefix-ultrametric} with \(\vartheta=1/3\) through \(\iota\) gives
an ultrametric \(d_C\) on that fractal.  If \(N\) is the first differing
digit, then
\begin{equation}
  \frac13 d_C(\omega,\omega')
  \leq \abs{\iota(\omega)-\iota(\omega')}
  \leq d_C(\omega,\omega').
  \label{eq:cantor-euclidean-comparison}
\end{equation}
Thus the familiar Euclidean geometry of the Cantor set is bi-Lipschitz
equivalent to an ultrametric geometry, although the restricted Euclidean
distance itself is not ultrametric.  This example is a useful warning that
ultrametricity depends on the chosen distance, not only on the topology.
\end{example}

\paragraph{Valuations produce arithmetic ultrametrics.}
The strong triangle inequality is the defining feature of non-Archimedean
absolute values.  Congruence at increasing precision plays exactly the role
of agreement along an increasingly long prefix.

\begin{example}[Finite congruence spaces and the \(p\)-adic integers]
\label{ex:padic-ultrametric}
Fix a prime \(p\).  For distinct residues
\(a,b\in\mathbb Z/p^L\mathbb Z\), let
\[
  n(a,b)=\max\{0\leq k<L:a\equiv b\pmod {p^k}\},
  \qquad d_p(a,b)=p^{-n(a,b)},
\]
and set \(d_p(a,a)=0\).  Congruence modulo \(p^k\) gives the nested
partitions, so this is a finite ultrametric.  Passing to the inverse limit
gives the compact ring of \(p\)-adic integers \(\mathbb Z_p\), with
\begin{equation}
  d_p(x,y)=\abs{x-y}_p=p^{-v_p(x-y)}.
  \label{eq:padic-ultrametric}
\end{equation}
The valuation inequality
\(v_p(x+y)\geq\min\{v_p(x),v_p(y)\}\) is precisely the ultrametric
inequality.  Balls are congruence classes modulo \(p^k\).  The field
\(\mathbb Q_p\) has the same local geometry but is not compact; compact
subsets such as \(\mathbb Z_p\), or probability measures supported on compact
balls, are therefore the direct setting for the compact GW statement.
\end{example}

\begin{example}[Formal series and non-Archimedean vector spaces]
\label{ex:valued-field-ultrametric}
For a finite field \(\mathbb F_q\), let
\(\operatorname{ord}_t(f)\) be the lowest nonzero exponent in a Laurent
series \(f\in\mathbb F_q((t))\).  Then
\[
  d(f,g)=q^{-\operatorname{ord}_t(f-g)}
\]
is ultrametric, and the unit ball \(\mathbb F_q[[t]]\) is compact.  More
generally, if \(\mathbb K\) is any non-Archimedean valued field, then
\(d(x,y)=\lVert x-y\rVert_\infty\) is ultrametric on
\(\mathbb K^m\), as well as on sequence spaces such as
\(c_0(\mathbb N,\mathbb K)\).  Hence ultrametric spaces are not confined to
trees: they also arise as finite- and infinite-dimensional algebraic state
spaces.
\end{example}

\paragraph{Closure rules generate more elaborate spaces.}
Several constructions preserve the strong triangle inequality and can be
used to build new GW input spaces from the preceding examples.

\begin{proposition}[Elementary closure properties]
\label{prop:ultrametric-closure}
Let \((\mathcal X,d)\) and \((\mathcal X_j,d_j)\), \(1\leq j\leq m\), be
ultrametric spaces.
\begin{enumerate}
  \item Every subset of \(\mathcal X\), with the restricted distance, is
  ultrametric.
  \item If \(\varphi:\mathbb R_+\to\mathbb R_+\) is nondecreasing,
  \(\varphi(0)=0\), and \(\varphi(r)>0\) for \(r>0\), then
  \(\varphi\circ d\) is ultrametric.  In particular, powers, truncations and
  bounded transforms such as \(d/(1+d)\) preserve ultrametricity.
  \item The max product
  \[
    d_{\max}(x,x')=\max_{1\leq j\leq m}d_j(x_j,x'_j)
  \]
  is ultrametric on \(\prod_j\mathcal X_j\).
  \item On the nonempty compact subsets \(A,B\subset\mathcal X\), the
  Hausdorff distance
  \begin{equation}
    d_H(A,B)=\max\left\{
      \sup_{x\in A}\inf_{y\in B}d(x,y),
      \sup_{y\in B}\inf_{x\in A}d(x,y)
    \right\}
    \label{eq:ultrametric-hausdorff}
  \end{equation}
  is ultrametric.  If \(\mathcal X\) is compact, then its Hausdorff
  hyperspace is compact as well.
\end{enumerate}
\end{proposition}

\begin{proof}
Restriction is immediate.  Monotonicity gives
\(\varphi(d(x,z))\leq
\varphi(\max\{d(x,y),d(y,z)\})\), and a nondecreasing function commutes with
the maximum of two scalars.  Taking the maximum of the componentwise strong
triangle inequalities proves the product claim.

For the Hausdorff claim, fix
\(r>\max\{d_H(A,B),d_H(B,C)\}\).  Every \(x\in A\) lies within \(r\) of some
\(y\in B\), which lies within \(r\) of some \(z\in C\).  Ultrametricity gives
\(d(x,z)<r\).  The same argument in the reverse direction proves
\(d_H(A,C)\leq\max\{d_H(A,B),d_H(B,C)\}\).  Compactness of the hyperspace
follows by approximating compact sets with subsets of a common finite
\(\varepsilon\)-net and using completeness.
\end{proof}

\begin{example}[Profinite groups and inverse limits]
\label{ex:profinite-ultrametric}
Let \(G=G_0\supset G_1\supset\cdots\) be finite-index normal subgroups with
\(\bigcap_kG_k=\{e\}\), and consider the inverse-limit completion.  For
\(0<\vartheta<1\), define
\[
  N(g,h)=\sup\{k:g^{-1}h\in G_k\},
  \qquad d(g,h)=\vartheta^{N(g,h)},
\]
with \(\vartheta^\infty=0\).  Two elements are close when they agree in many
finite quotients.  Since agreement through level \(k\) is an equivalence
relation, this is a compact ultrametric group; normality makes the distance
bi-invariant.  The additive group of \(\mathbb Z_p\) is recovered from
\(G_k=p^k\mathbb Z_p\), while products of finite groups give prefix metrics
on increasingly refined discrete coordinates.
\end{example}

\begin{example}[The local topology on rooted graphs]
\label{ex:rooted-graph-ultrametric}
Let \(\mathscr G_D\) be the isomorphism classes of connected rooted simple
graphs with degree at most \(D\).  For rooted graphs \((G,o)\) and
\((G',o')\), let
\(R(G,G')\) be the largest radius for which the rooted radius-\(R\)
neighborhoods are isomorphic, and set
\begin{equation}
  d_{\mathrm{loc}}(G,G')=2^{-R(G,G')},
  \qquad 2^{-\infty}=0.
  \label{eq:rooted-graph-ultrametric}
\end{equation}
Agreement of two rooted balls through radius \(R\) is transitive, so
\eqref{eq:rooted-graph-ultrametric} is ultrametric.  Bounded degree leaves
only finitely many rooted ball types at each radius; a diagonal argument then
shows that \(\mathscr G_D\) is compact.  Weak convergence of probability laws
on this space is the local weak, or Benjamini--Schramm, convergence used for
large random graphs \citep{BenjaminiSchramm2001}.
\end{example}

\begin{example}[Compact-set-valued data]
\label{ex:hyperspace-ultrametric}
Starting from any compact example above, one may take its nonempty compact
subsets and compare them with \eqref{eq:ultrametric-hausdorff}.  Iterating the
construction gives spaces of hierarchical set-valued objects, all still
compact and ultrametric.  For a finite dendrogram, a set is observed at
resolution \(r\) only through the collection of level-\(r\) clusters it
meets; two sets are within \(r\) exactly when these collections agree.  This
makes the threshold partition behind the induced Hausdorff ultrametric
explicit.
\end{example}

\begin{example}[Random tree boundaries and hierarchical cascades]
\label{ex:random-cascade-ultrametric}
Let \(\mathcal T\) be an infinite locally finite rooted tree and let
\(\partial\mathcal T\) be its infinite rays.  Choose radii
\(r_0>r_1>\cdots\downarrow0\).  If \(x\wedge y\) is the deepest common vertex
of two rays, then
\[
  d(x,y)=r_{\lvert x\wedge y\rvert}\quad(x\neq y),
  \qquad d(x,x)=0,
\]
is an ultrametric, and \(\partial\mathcal T\) is compact.  Assigning random
child weights that sum to one at each vertex produces a random probability
measure whose cylinder mass is the product of the weights along its branch.
The geometry remains ultrametric while the measure becomes highly
inhomogeneous.  Such random hierarchical measures are the elementary
template behind multiplicative cascades and the ultrametric organization of
states in spin-glass models; Ruelle probability cascades provide a canonical
infinite-level realization \citep{RammalToulouseVirasoro1986,Ruelle1987}.
\end{example}

\begin{remark}[Which examples enter the GW theorem]
Every finite example above satisfies \cref{cor:ultrametric-extreme-minimizer}.
The finite-alphabet sequence spaces, Cantor set, \(\mathbb Z_p\), compact
unit balls such as \(\mathbb F_q[[t]]\), profinite groups, bounded-degree
rooted graph space, compact Hausdorff hyperspaces and locally finite tree
boundaries all satisfy the hypotheses of \cref{thm:ultrametric-concavity}.
Noncompact examples such as Baire space, \(\mathbb Q_p\), full Laurent-series
fields and infinite-dimensional sequence spaces require either compact
restrictions or an extension with moment and integrability assumptions.
Ordinary shortest-path tree metrics do not enter this list unless their
restriction has separately been shown to be ultrametric.
\end{remark}

\section{Compact ultrametric spaces}
\label{sec:ultrametric}

Ultrametric geometry supplies a concrete and useful class where distance
thresholds are automatically conditionally negative.  At each scale, the
space breaks into disjoint clusters, turning the key product quadratic form
into a sum of squares.

\begin{theorem}[Conditional concavity on compact ultrametric spaces]
\label{thm:ultrametric-concavity}
Let \((\mathcal X,d_{\mathcal X})\) and
\((\mathcal Y,d_{\mathcal Y})\) be compact ultrametric spaces.  For arbitrary
probability measures \(\alpha\in\Pcal(\mathcal X)\) and
\(\beta\in\Pcal(\mathcal Y)\), and every submodular mixed-measure distortion
\(\ell\), the objective
\[
  \pi\longmapsto
  \iint
  \ell\bigl(d_{\mathcal X}(x,x'),d_{\mathcal Y}(y,y')\bigr)
  \dd\pi(x,y)\dd\pi(x',y')
\]
is concave on \(\Piopt(\alpha,\beta)\).  In particular, this holds for every
power distortion \(\ell(a,b)=\abs{a-b}^p\), \(p\geq1\).
\end{theorem}

\begin{proof}
Fix \(r>0\).  In an ultrametric space, the relation
\(x\sim_r x'\) if \(d_{\mathcal X}(x,x')<r\) is an equivalence relation.
Compactness implies that its partition \(\mathscr P_r^{\mathcal X}\) has
finitely many classes.  Therefore
\[
  H_r^{d_{\mathcal X}}(x,x')
  =1-\sum_{C\in\mathscr P_r^{\mathcal X}}
  \mathbf 1_C(x)\mathbf 1_C(x').
\]
For every finite signed measure \(\sigma\) of total mass zero,
\[
  \iint H_r^{d_{\mathcal X}}(x,x')\dd\sigma(x)\dd\sigma(x')
  =-\sum_{C\in\mathscr P_r^{\mathcal X}}\sigma(C)^2\leq0.
\]
Thus every threshold kernel is CND, in agreement with
\cref{prop:threshold-ultrametric}, and the same argument applies on
\(\mathcal Y\).

More directly, if \(\xi\in\Tcal_{\mathcal X,\mathcal Y}\), then the analogue
of \eqref{eq:threshold-product} is
\begin{equation}
  I_\xi(r,s)
  =
  \sum_{C\in\mathscr P_r^{\mathcal X}}
  \sum_{D\in\mathscr P_s^{\mathcal Y}}
  \xi(C\times D)^2
  \geq0.
  \label{eq:ultrametric-squares}
\end{equation}
Indeed, the constant term and the terms involving only one partition vanish
because \(\xi\) has zero total mass and zero marginals.  The same Fubini
calculation as in \eqref{eq:mixed-measure-curvature} gives
\[
  \iint \ell\bigl(d_{\mathcal X}(x,x'),d_{\mathcal Y}(y,y')\bigr)
  \dd\xi(x,y)\dd\xi(x',y')
  =-\int I_\xi(r,s)\dd\omega_\ell(r,s)\leq0.
\]
The quadratic segment expansion used in \cref{prop:curvature-criterion}
therefore proves concavity.  The power case follows from
\cref{cor:power-distortions}.
\end{proof}

\begin{remark}[The endpoint \(p=\infty\)]
The usual \(p=\infty\) GW distortion is an essential supremum rather than a
quadratic functional of the coupling.  The direction-space notion of
conditional concavity used here therefore does not pass to that endpoint.
\end{remark}

\section{Conclusion and recommended book statement}
\label{sec:conclusion}

The clean presentation should begin with the master identity
\eqref{eq:mixed-measure-curvature}, not with two apparently unrelated proofs.
For genuine metrics, it yields a hierarchy of certificates:
\[
  \text{ultrametric}
  \quad\Longleftrightarrow\quad
  \text{threshold-CND}
  \quad\Longrightarrow\quad
  \text{CND}.
\]
At \(p=2\), the constant scale weight allows one to use only the weakest,
integrated CND/CPD certificate, as formalized in
\cref{cor:quadratic-integrated-certificate}.  Under the stronger ultrametric
certificate, the same formula controls every pair of scales and therefore
every \(p\geq1\), as well as every submodular mixed-measure distortion.  Thus
the ultrametric theorem genuinely contains the quadratic case; what is special
about \(p=2\) is that its hypothesis can be weakened.

The counterexamples in \cref{thm:no-cnd-extension} show that this weakening
cannot simply be reused for \(p\neq2\).  In finite ultrametric spaces,
\cref{cor:ultrametric-extreme-minimizer} additionally yields sparse and, for
uniform marginals, permutation optimizers.  This organization also prevents
confusing the outer GW exponent with powers used inside the base
dissimilarities.

\bibliographystyle{plainnat}
\bibliography{references}

\end{document}
